\section{Chapter 5}

\begin{quotex}
\begin{verse}
The Universal Heaven and Earth is not humane compassionate \\
All things (products) use them as means\\
True Men are not humane (\emph{jen})\\
Beings use them as means\\
(Like) emptiness between Heaven and Earth in its Virtue\\
The Principle is like a bellows\\
It is emptied yet unexhausted\\
Even if creating inexhaustibly with its motion\\
(Thus) it is vain to multiply words and projects (in the care for individuals)\\
Keeping in the middle is the best.
\end{verse}
\end{quotex}

In the word translated as ``means'' (second and fourth lines) is literally ``straw dogs'' (\emph{ch'a kou}), likenesses fabricated in ancient China for certain rites and then abandoned, thrown away. It wants to express that the Absolute considers individuals only as a universal function, i.e., as a function of that which transcends them, for every other verse, that is for what are simple ephemeral apparitions in the current of forms, the Principle not taking care of them (non-humanizing), against the religious conception of Providence and the God of Love.

Since he reproduces in himself the detached quality of the Principle, the True Man likewise ignores human concern, even as the Confucian \emph{jen} (sympathy, sociality), from some Taoist comments called the ``small virtue'' or the ``lower virtue''. In the last two lines, in word and action in view of the particular as such, the action from the ``center'' is counterposed, that is to say, from a neutrality free from preferences and sentimentality, in view not of part, but of the whole, of the Universal. It makes the virtue of the Tao from the model, which even if producing inexhaustibly remains ``empty'' (image of the bellows), that is, pure simple, ``free from the I''. See Chuang Tzu (V,2) who says, referring to Wang t'ai, whose influences ``derive from its transcendence'': Reaching perfect impassability ``lf beings evolve according to their destinies and maintains himself as the immobile center of all fates'' (as in another passage, II, e, it is said, of the Principle: ``The immobile center of a circle on the circumference of which run all contingencies, distinctions, and individuals'' since ``there are two ways to consider beings, as distinct individuals or as all in a great whole''. It is the imperturbability of a spirit that, ``higher than heaven, earth, and all beings, lives in a body. Absolutely independent, he is lord of men''.

There are those who have wanted to see the reflection of a cold Machiavellianism without scruples in the ideas just expressed. Also the descent down to the political plane that could be right only if the reference point was a powerful individual's will to dominate, that is, the intensification of a particular I who, on the contrary, is the first thing that True Men destroy in themselves.

About the last line, we read also in the Chung Yung: ``To maintain the middle path is perfection, but there are very few in the world who can keep it''. That is to say, the metaphysical meaning is opposed to that of moderation, of the practical ``golden mean''. About the noted relativity of the divergence of the views of Las Tzu and Confucius, we can cite, still in the same context, the following passage of the latter: ``Who rules according to the virtue of Heaven resembles the polar star: it is immobile, but all revolve orderly around it'' (\emph{Lun you}, II, 1 Analects ).

\textbf{Chinese text and literal translation}


\paragraph{Chapter 5 (\begin{CJK*}{UTF8}{bsmi}第五章\end{CJK*})}
\columnratio{.4}
\begin{paracol}{2}
\begin{verse}
\begin{CJK*}{UTF8}{bsmi}
天地不仁\\
以萬物為芻狗；\\
聖人不仁，\\
以百姓為芻狗。\\
天地之閒其猶橐籥乎？\\
虛而不屈，\\
動而愈出。\\
多言數窮，\\
不如守中。
\end{CJK*}
\end{verse}
\switchcolumn
\begin{verse}
The sky and the earth do not care,\\
They regard the myriad things as straw dogs;\\
The sage does not care,\\
He regards people as straw dogs.\\
The space between the sky and the earth, how much is it like large bellows!\\
Empty but endless,\\
Just move and wind will be produced;\\
Much talk soon comes to nothing.\\
It is better to be in between extremes.
\end{verse}
\end{paracol}


\flrightit{Posted on 2020-08-22 by Lao Tzu}

\begin{center}* * *\end{center}

\begin{footnotesize}
\begin{sffamily}
\texttt{Sensitive Ears on 2020-08-27 at 01:19 said:}

These posts would be improved by audio recordings of the Chinese.

\hfill

\texttt{Cologero on 2020-08-27 at 06:19 said:}

Mr. Ears, you remind me of a bad girlfriend ... whatever I do is never good enough.

Do you prefer Mandarin or Cantonese for the recording?


\hfill

\end{sffamily}
\end{footnotesize}

